Figure~\ref{fig:acc} shows accuracy against $h$ for each feature strength; Table~\ref{tab:main} gives the $\mu=1$ cells with seed standard deviations; Tables~\ref{tab:win} and~\ref{tab:gap} give the winner and the GCN$-$MLP gap in all 33 cells. The MLP row is constant in $h$ because, per seed, only the edges change.

\begin{figure*}[t]\centering\includegraphics[width=\textwidth]{acc_vs_h.pdf}
\caption{Test accuracy versus edge homophily for the four systems at three feature strengths (mean over 5 seeds, bars show the standard error). The dotted line marks $h=1/K$.}\label{fig:acc}\end{figure*}

\begin{table}[t]\centering\caption{Test accuracy at $\mu=1$ (mean $\pm$ std over 5 seeds; bold = best mean); last column: paired $p$ for GCN vs MLP.}\label{tab:main}
\resizebox{\columnwidth}{!}{\begin{tabular}{rccccc}
\toprule
$h$ & MLP & GCN & H2GCN & LP & $p_{\rm GCN-MLP}$ \\
\midrule
0.0 & 0.539$\pm$0.029 & 0.709$\pm$0.039 & \textbf{0.911}$\pm$0.025 & 0.321$\pm$0.022 & \rhval{cmp/main/mlp/h0.0_mu1.0/accuracy/paired_p} \\
0.1 & 0.539$\pm$0.029 & 0.542$\pm$0.034 & \textbf{0.724}$\pm$0.037 & 0.327$\pm$0.008 & \rhval{cmp/main/mlp/h0.1_mu1.0/accuracy/paired_p} \\
0.2 & 0.539$\pm$0.029 & 0.444$\pm$0.020 & \textbf{0.587}$\pm$0.023 & 0.311$\pm$0.038 & \rhval{cmp/main/mlp/h0.2_mu1.0/accuracy/paired_p} \\
0.3 & \textbf{0.539}$\pm$0.029 & 0.449$\pm$0.023 & 0.535$\pm$0.039 & 0.312$\pm$0.027 & \rhval{cmp/main/mlp/h0.3_mu1.0/accuracy/paired_p} \\
0.4 & \textbf{0.539}$\pm$0.029 & 0.506$\pm$0.023 & 0.536$\pm$0.040 & 0.389$\pm$0.023 & \rhval{cmp/main/mlp/h0.4_mu1.0/accuracy/paired_p} \\
0.5 & 0.539$\pm$0.029 & \textbf{0.638}$\pm$0.031 & 0.635$\pm$0.034 & 0.501$\pm$0.022 & \rhval{cmp/main/mlp/h0.5_mu1.0/accuracy/paired_p} \\
0.6 & 0.539$\pm$0.029 & \textbf{0.792}$\pm$0.024 & 0.770$\pm$0.026 & 0.664$\pm$0.016 & \rhval{cmp/main/mlp/h0.6_mu1.0/accuracy/paired_p} \\
0.7 & 0.539$\pm$0.029 & \textbf{0.904}$\pm$0.011 & 0.890$\pm$0.018 & 0.838$\pm$0.023 & \rhval{cmp/main/mlp/h0.7_mu1.0/accuracy/paired_p} \\
0.8 & 0.539$\pm$0.029 & \textbf{0.970}$\pm$0.010 & 0.967$\pm$0.010 & 0.960$\pm$0.014 & \rhval{cmp/main/mlp/h0.8_mu1.0/accuracy/paired_p} \\
0.9 & 0.539$\pm$0.029 & 0.992$\pm$0.007 & 0.994$\pm$0.003 & \textbf{0.996}$\pm$0.003 & \rhval{cmp/main/mlp/h0.9_mu1.0/accuracy/paired_p} \\
1.0 & 0.539$\pm$0.029 & 0.997$\pm$0.003 & 0.998$\pm$0.002 & \textbf{1.000}$\pm$0.000 & \rhval{cmp/main/mlp/h1.0_mu1.0/accuracy/paired_p} \\
\bottomrule
\end{tabular}}\end{table}

\begin{table*}[t]\centering\caption{Best system (by mean over 5 seeds) in each of the 33 cells.}\label{tab:win}
\begin{tabular}{lccccccccccc}
\toprule
$\mu$ \textbackslash\ $h$ & 0.0 & 0.1 & 0.2 & 0.3 & 0.4 & 0.5 & 0.6 & 0.7 & 0.8 & 0.9 & 1.0 \\
\midrule
0.5 & H2GCN & H2GCN & H2GCN & H2GCN & GCN & GCN & LP & LP & LP & LP & LP \\
1 & H2GCN & H2GCN & H2GCN & MLP & MLP & GCN & GCN & GCN & GCN & LP & LP \\
2 & H2GCN & H2GCN & H2GCN & MLP & MLP & H2GCN & H2GCN & H2GCN & GCN & H2GCN & LP \\
\bottomrule
\end{tabular}\end{table*}

\begin{table*}[t]\centering\caption{GCN minus MLP test accuracy in percentage points (mean over 5 seeds); negative values are cells where GCN is below the MLP.}\label{tab:gap}
\begin{tabular}{lccccccccccc}
\toprule
$\mu$ \textbackslash\ $h$ & 0.0 & 0.1 & 0.2 & 0.3 & 0.4 & 0.5 & 0.6 & 0.7 & 0.8 & 0.9 & 1.0 \\
\midrule
0.5 & +13.9 & +0.7 & -3.2 & -1.9 & +1.0 & +11.6 & +23.4 & +36.9 & +49.0 & +56.1 & +59.1 \\
1 & +17.0 & +0.3 & -9.5 & -9.1 & -3.3 & +9.9 & +25.3 & +36.5 & +43.1 & +45.3 & +45.8 \\
2 & +7.1 & -10.0 & -22.7 & -25.3 & -16.4 & -4.3 & +6.7 & +13.3 & +16.4 & +16.6 & +16.9 \\
\bottomrule
\end{tabular}\end{table*}

\textbf{H1: GCN below the MLP at mid homophily (partly supported, not where we registered it).} At $\mu=1$, GCN is below the MLP at $h=0.3$ (0.449 vs 0.539, $p=0.007$) and at $h=0.4$ the mean is also lower (0.506) but the paired test is not significant ($p=0.142$); at $h=0.5$ GCN is \emph{above} the MLP (0.638, $p=0.014$). By the registered criterion (negative and $p<0.05$ in the cells $0.3,0.4,0.5$) only one of three cells qualifies, so H1 is supported at $h=0.3$ and not as a statement about the whole band. The deficit is in fact larger below the registered cells: Table~\ref{tab:gap} gives $-9.5$ points at $h=0.2$ ($p=0.001$ in Table~\ref{tab:main}) and $-9.1$ at $h=0.3$. It grows with feature strength, from $-3.2$ points at $\mu=0.5$ to $-22.7$ and $-25.3$ at $\mu=2$ ($h=0.2,0.3$), consistent with the picture that mixing in neighbours of other classes destroys good features; we did not run an experiment isolating that mechanism. For $\mu=0.5$ the dip is small ($-3.2$ and $-1.9$ points at $h=0.2$ and $0.3$) and we regard it as weak evidence. GCN also beats the MLP at $h=0$ (+17.0 at $\mu=1$): with three classes, every neighbour of a node is from a different class, so aggregated neighbourhoods remain class-dependent. This U-shape is the shape suggested by~\cite{ma2021is}, but our model does not test their neighbourhood-distribution explanation.

\textbf{H2: ego/neighbour separation removes the dip (refuted as stated).} H2GCN-style is the best system in all cells with $h\le0.2$ (Table~\ref{tab:win}) and is far ahead under heterophily: at $h=0$, $\mu=1$ it reaches 0.911 versus 0.709 for GCN and 0.539 for MLP. At $h=0.3$--$0.4$, $\mu=1$ it recovers the MLP level (0.535, 0.536 vs 0.539) but does not exceed it. The registered criterion, within 0.02 of $\max$(MLP, GCN) in every cell, fails in 6 of 33 cells (Table~\ref{tab:viol}), the largest at $\mu=0.5$, $h=0.6$ where it trails GCN by 0.078, and one with strong features ($\mu=2$, $h=0.3$) where it trails the MLP by 0.027. So the separation reduces but does not remove the mid-homophily dip, and at moderate-to-high homophily the extra neighbour weights cost accuracy when training data are scarce relative to parameters (we did not isolate this).

\begin{table}[t]\centering\caption{Cells where H2GCN-style is more than 0.02 below the better of MLP and GCN (test accuracy, mean over 5 seeds).}\label{tab:viol}
\resizebox{\columnwidth}{!}{\begin{tabular}{rrccc}
\toprule
$\mu$ & $h$ & MLP & GCN & H2GCN \\
\midrule
0.5 & 0.4 & 0.390 & 0.400 & 0.373 \\
0.5 & 0.5 & 0.390 & 0.506 & 0.454 \\
0.5 & 0.6 & 0.390 & 0.624 & 0.546 \\
0.5 & 0.7 & 0.390 & 0.759 & 0.731 \\
1 & 0.6 & 0.539 & 0.792 & 0.770 \\
2 & 0.3 & 0.830 & 0.577 & 0.804 \\
\bottomrule
\end{tabular}}\end{table}

\textbf{H3: label propagation (refuted).} LP is the best system only in the high-homophily corner (Table~\ref{tab:win}): $h\ge0.6$ at $\mu=0.5$, $h\ge0.9$ at $\mu=1$ and $h=1$ at $\mu=2$. The first part of H3 (best only with weak features and $h\ge0.8$) is false because LP also wins at $\mu=1,2$ in the extreme cells, although there its margin over GCN/H2GCN is tiny (e.g.\ 1.000 vs 0.997 and 0.998 at $h=1$, $\mu=1$; at $h=1$ LP's $\pm0.000$ is not a bug: with $h=1$ the graph splits into class-pure components that the 20\% training labels cover). It is below the MLP at $h\le0.4$, $\mu=1$ (0.389 at $h=0.4$ vs 0.539).

\textbf{H4: crossover moves with informativeness (supported on the homophilous branch).} Reading the crossover as the lowest $h\ge1/3$ at which GCN's mean exceeds the MLP's, it is $h=0.4$ for $\mu=0.5$, $0.5$ for $\mu=1$ and $0.6$ for $\mu=2$ (Table~\ref{tab:gap}): the stronger the features, the more homophily GCN needs. At $\mu=0.5$ and $h=0.4$ the gap is $+1.0$ points, inside seed noise, so the $\mu=0.5$ crossover is uncertain by about one grid step, and the ordering has only three levels of $\mu$. This is a descriptive statement at 0.1 resolution, and as noted in Section~\ref{sec:setup} it uses a post-hoc restriction of the registered wording.

\textbf{Where the methods fail.} GCN: $h\in[0.2,0.4]$ for $\mu\ge1$. H2GCN-style: moderate homophily with weak features (Table~\ref{tab:viol}). LP: any $h\le0.5$ at $\mu\ge1$, and it uses no features. MLP: wins only at $h=0.3$--$0.4$, $\mu\ge1$ (Table~\ref{tab:win}), the band where the graph is almost uninformative; its win is by 0.4 points over H2GCN-style at $\mu=1$ and $h=0.3$, which is within noise.
